\section{Related Work}
\label{sec:related}

\subsection{Autoresearch and LLM-Driven Optimization}

\citet{karpathy2026autoresearch} introduced the paradigmatic autoresearch loop
for neural network hyperparameter search: an LLM reads a training script,
proposes a configuration change, executes training for a fixed budget, measures
validation loss, and accepts or rejects the change.
Iterated, this constitutes a form of LLM-guided hill climbing in configuration
space, where the LLM's world knowledge serves as an implicit prior over
promising changes and training outcomes provide gradient-free feedback.

AutoResearchClaw~\citep{aiminglab2026claw} extends this framework with
multi-batch parallelism: several candidate configurations are evaluated
simultaneously, and the best is promoted.
This increases the effective branching factor of search without altering the
underlying acceptance mechanism.

EvoScientist~\citep{evoscientist2026} introduces persistent experience memory:
lessons from prior runs are summarized and injected into future proposals,
enabling cross-run learning.
Both of these enhancements were designed by human researchers who inspected the
prior system's code and identified architectural gaps.
In all three systems, the structural decisions---when to accept, how to
propose, what state to maintain---are made by human designers, not by the
system itself.

\subsection{Bilevel Optimization}

Bilevel optimization~\citep{colson2007overview,sinha2018review} studies
problems of the form
$\min_{\phi}\, F\!\left(\phi,\, \theta^{*}(\phi)\right)$
subject to $\theta^{*}(\phi) \in \arg\min_{\theta}\, f(\theta, \phi)$,
where an upper-level objective $F$ depends on the optimal solution of a
lower-level problem parameterized by $\phi$.
Applications include meta-learning~\citep{franceschi2018bilevel},
neural architecture search~\citep{liu2018darts}, and hyperparameter
optimization~\citep{feurer2019hyperparameter}.
In our setting the upper level optimizes the search \emph{mechanism} $\phi$
(the runner code) and the lower level optimizes the task performance $\theta$
(the training configuration).
The key departure from classical bilevel optimization is that $\phi$ is a
program---a discrete artifact produced by code generation---rather than a
real-valued parameter vector.

\subsection{LLM-Based Code Generation for Research}

AlphaCode~\citep{li2022competition} and Codex~\citep{chen2021evaluating}
demonstrated that LLMs can write functionally correct programs from natural
language specifications.
FunSearch~\citep{romera2024mathematical} extended this to scientific
discovery, using an LLM to iteratively generate and evaluate mathematical
programs, finding new results in combinatorics.
Most directly related is the line of work on LLM-driven algorithm
design~\citep{liu2024evolution,lehman2023evolution}, in which LLMs propose
novel algorithmic variants that are then evaluated on benchmark tasks.
Our Level~2 agent applies the same code generation capacity to a different
target: rather than generating task-level programs, it generates
\emph{search mechanism code} that is injected into the inner loop at runtime.

\subsection{Meta-Learning and Algorithm Configuration}

Meta-learning~\citep{hospedales2021meta} trains models to learn efficiently
from few examples by optimizing across a distribution of tasks.
Algorithm configuration~\citep{hutter2011smac} and algorithm
selection~\citep{rice1976algorithm} choose among candidate algorithms or
parameter settings for a given problem instance.
Portfolio methods~\citep{xu2008satzilla} maintain a library of algorithms and
select among them.
\methodname{} operates in a similar spirit---the outer loop selects or
generates a search mechanism---but uses LLM code generation rather than
gradient-based meta-optimization or a fixed portfolio.

\subsection{Position of This Work}

The key distinction between \methodname{} and all prior work is the
\emph{target} of the outer loop.
Level~1.5 is closest to existing outer loops (curriculum schedulers, adaptive
configuration) in that it adjusts parameters of the existing mechanism.
Level~2 is categorically different: it generates code that \emph{replaces} the
mechanism entirely.
To our knowledge, \methodname{} is the first system in which an autonomous
outer loop writes and injects code that modifies the structural logic of the
inner autoresearch loop at runtime, using the same model that runs the inner
loop.
